%=============================================================================
% The χ Field of CP²×S³: Cosmological Completion
% of Density Field Dynamics
%
% Companion paper to DFD Unified Review v3.3
% Revised with honest grading and expanded derivations
%=============================================================================
\documentclass[aps,prd,twocolumn,superscriptaddress,showpacs,floatfix,nofootinbib]{revtex4-2}

% Mathematics
\usepackage{amsmath,amssymb,amsfonts,amsthm}
\usepackage{bm}
\usepackage{mathtools}
\usepackage[utf8]{inputenc}

% Graphics
\usepackage{graphicx}
\usepackage{xcolor}

% Tables
\usepackage{booktabs}
\usepackage{array}
\usepackage{multirow}

% References
\usepackage[colorlinks=true,linkcolor=blue!60!black,citecolor=blue!60!black,
            urlcolor=blue!60!black]{hyperref}

% Theorem environments
\newtheorem{theorem}{Theorem}
\newtheorem{lemma}[theorem]{Lemma}
\newtheorem{corollary}[theorem]{Corollary}
\newtheorem{proposition}[theorem]{Proposition}
\theoremstyle{definition}
\newtheorem{definition}[theorem]{Definition}
\theoremstyle{remark}
\newtheorem*{remark}{Remark}

% Custom commands
\newcommand{\CP}{\mathbb{CP}}
\newcommand{\Sph}{S}
\newcommand{\MP}{M_{\mathrm{P}}}
\newcommand{\MPbar}{\bar{M}_{\mathrm{P}}}
\newcommand{\keV}{\,\mathrm{keV}}
\newcommand{\GeV}{\,\mathrm{GeV}}
\newcommand{\eV}{\,\mathrm{eV}}
\newcommand{\MeV}{\,\mathrm{MeV}}
\newcommand{\Mpc}{\,\mathrm{Mpc}}
\newcommand{\Tr}{\operatorname{Tr}}
\newcommand{\astar}{a_*}
\newcommand{\azero}{a_0}
\newcommand{\kmax}{k_{\max}}
\newcommand{\Ngen}{N_{\mathrm{gen}}}
\newcommand{\LCDM}{\Lambda\mathrm{CDM}}

% Grading macro for honest assessment
\newcommand{\grade}[1]{\textsf{\footnotesize [#1]}}

\begin{document}

\title{The $\chi$~Field of $\CP^2\times\Sph^3$:\\
Cosmological Completion of Density Field Dynamics}

\author{Gary Alcock}
\affiliation{Independent Researcher}
\email{gary@densityfielddynamics.com}

\date{\today}

\begin{abstract}
Density Field Dynamics (DFD) explains galactic and cluster
dynamics through the scalar refractive field~$\psi$, without
invoking dark matter particles. However, we prove that $\psi$
alone cannot reproduce the observed matter power spectrum:
the MOND-enhanced gravitational potential is an odd function
of the baryon perturbation and therefore averages to zero
over pre-recombination acoustic oscillations, providing no
template for structure formation. The same $\CP^2\times\Sph^3$
topology that produces $\psi$ mandates, through the K\"unneth
theorem ($b_3=1$), a second scalar~$\chi$: the Chern--Simons
period on $\Sph^3$. We perform the Kaluza--Klein reduction
to four dimensions, showing that $\chi$ acquires a kinetic
term from the harmonic $3$-form normalization and a potential
from the Chern--Simons partition function. The resulting
dynamics are pressure-free ($w=0$, $c_s^2=0$) with no
photon coupling---identical to cold dark matter at
cosmological scales. The Toeplitz determinant-scaling
mechanism (already proved for $H_0$ and $M_R$ in the
unified review) yields a mass
$m_\chi \approx 96\keV$ when applied to the
$8$-dimensional automorphism space of $\CP^2$ and
the $5$-dimensional chiral multiplet space. The spectral-trace
factorization on the product geometry gives an abundance ratio
$\Omega_\chi/\Omega_b = 16/3 = 5.333$, consistent with
Planck~2018 ($5.36\pm 0.07$) and excluding the Standard
Model without right-handed neutrinos at $5.6\sigma$. The
$\Sph^3$ orientation symmetry ($\mathbb{Z}_2$: $\chi\to-\chi$)
forbids all photon couplings, ensuring absolute stability.
We present each result with an explicit assessment of its
derivation status, distinguishing theorem-grade, derived,
and conjectured claims. Testable predictions are identified
for Simons Observatory, ELT/ANDES, and current Planck data.
\end{abstract}

\maketitle

%==========================================================================
\section{Introduction and Motivation}
\label{sec:intro}
%==========================================================================

Density Field Dynamics (DFD)~\cite{Alcock2025} replaces
General Relativity's curved-spacetime ontology with a scalar
refractive field~$\psi$ on flat~$\mathbb{R}^3$, governed by
\begin{equation}
\label{eq:psi-field}
\nabla\cdot\bigl[\mu(|\nabla\psi|/\astar)\,\nabla\psi\bigr]
= -\frac{8\pi G}{c^2}\,\rho_b\,,
\end{equation}
where $\mu(x)=x/(1+x)$ is derived from the $\Sph^3$
composition law and $\astar = 2\azero/c^2$. The theory
reproduces all classical tests of General Relativity
(PPN parameters $\gamma=\beta=1$), explains galaxy rotation
curves without dark matter (175~SPARC galaxies), derives
$\alpha^{-1}=137.036$ from Chern--Simons theory on $\Sph^3$,
and accounts for apparent cosmic acceleration through the
$\psi$-screen optical effect~\cite{Alcock2025}.

The internal manifold $M_7=\CP^2\times\Sph^3$ is uniquely
selected by six vacuum axioms~\cite{Alcock2025b}. This
topology determines the Standard Model gauge group, three
generations ($\Ngen=\chi(\CP^2)=3$), and the strong-CP
solution $\bar\theta=0$ to all loop orders.

DFD~v3.3 explicitly identified the matter power
spectrum~$P(k)$ as a ``program item''---an open sector
where the theory had not yet been tested against
observations. The present work addresses this gap by
completing the zero-mode census of $\CP^2\times\Sph^3$
and identifying the $\chi$~field from $b_3=1$ as the
pre-recombination clustering component.

We emphasize that $\chi$ is \emph{not} ``dark matter''
in the conventional sense---an unknown particle postulated
to explain gravitational anomalies. It is a specific,
derived consequence of the topology that already produces
$\alpha$, the fermion masses, and the gravitational
field~$\psi$. DFD does not \emph{need} dark matter for
galactic dynamics (handled by $\psi$) or cosmic
acceleration (handled by the $\psi$-screen). The $\chi$~field
fills a complementary role: providing the pre-recombination
gravitational scaffolding that the $\psi$~field provably
cannot supply.

\subsection{Conventions and notation}

Throughout, $\MP = \sqrt{\hbar c/G} = 1.221\times 10^{19}\GeV$
is the full Planck mass and
$\MPbar = \MP/\sqrt{8\pi} = 2.435\times 10^{18}\GeV$ the
reduced Planck mass. DFD~v3.3 uses $\MP$ uniformly; the
$\chi$~sector naturally employs $\MPbar$ from the
Einstein--Hilbert normalization. We indicate which convention
is used in each formula. Numerical values are verified
against $\MP\alpha^n$ for the established $\alpha$-tower
rungs~\cite{Alcock2025}.

\subsection{Honest grading}

Each major result is accompanied by a derivation-status
label:
\begin{itemize}
\item \grade{Theorem}: proved from mathematical axioms
  with no physical assumptions beyond the DFD postulates.
\item \grade{Derived}: follows from a specific chain of
  reasoning with clearly identified assumptions.
\item \grade{Conjectured}: supported by strong evidence
  (numerical agreement, structural parallels) but with
  identified gaps in the derivation.
\end{itemize}
This grading system reflects the maturity of the
theoretical program, not uncertainty about the physics.
Results labeled \grade{Conjectured} identify specific
calculations that would elevate them to \grade{Derived}
or \grade{Theorem}.

%==========================================================================
\section{Why $\psi$ Alone Cannot Produce $P(k)$}
\label{sec:impossibility}
%==========================================================================

\begin{theorem}[Odd-Function Structural Obstruction \grade{Theorem}]
\label{thm:odd}
Let $\Phi_N$ denote the Newtonian gravitational potential from
baryonic perturbations, and define the DFD-enhanced potential
$\Phi_{\mathrm{DFD}} = \nu(|\Phi_N|)\,\Phi_N$ where
$\nu(y)=\tfrac{1}{2}(1{+}\sqrt{1{+}4/y})$. Then
$\langle\Phi_{\mathrm{DFD}}\rangle_t = 0$ over any
symmetric oscillation of~$\Phi_N$.
\end{theorem}

\begin{proof}
(i)~$\nu(y)$ depends on $|\Phi_N|$, so $\nu$ is even:
$\nu(|-\Phi_N|) = \nu(|\Phi_N|)$.
(ii)~Therefore $\Phi_{\mathrm{DFD}}(-\Phi_N)
= \nu(|\Phi_N|)(-\Phi_N) = -\Phi_{\mathrm{DFD}}(\Phi_N)$;
the enhanced potential is \emph{odd} in $\Phi_N$.
(iii)~Before recombination, baryons are tightly coupled to
photons and oscillate acoustically:
$\delta_b\propto\cos(kc_s\eta)$. The Poisson equation gives
$\Phi_N\propto\delta_b$, so $\Phi_N$ oscillates symmetrically.
(iv)~For any odd function $f$ and symmetric distribution~$P$,
$\langle f\rangle = \int f(x)P(x)\,dx = 0$.
\end{proof}

The theorem holds exactly in the linear acoustic regime.
Nonlinear corrections (mode coupling, second-order effects)
are suppressed by $\delta_b^2\sim 10^{-10}$ at
recombination and cannot overcome the structural zero.
A systematic computational campaign (over 120~independent
analyses, documented in
Ref.~\cite{Alcock2025_supp}, covering mode coupling,
$\sigma_\nabla$ self-regulation, temporal inertia, and
pre-recombination MOND enhancement) verified that no
nonlinear, non-acoustic, or mode-coupling mechanism
overcomes this obstruction.

\emph{Contrast with CDM.} Cold dark matter produces
$\langle\Phi_{\mathrm{CDM}}\rangle\neq 0$ because
$\delta_{\mathrm{CDM}}$ grows \emph{monotonically}---it is
decoupled from photons and has no acoustic oscillation.
The structural difference is not quantitative enhancement
but the \emph{sign symmetry} of the source: CDM breaks
it, $\psi$ preserves it.

Four additional arguments reinforce the theorem:

\emph{(A)~Elliptic constraint \grade{Theorem}.} The spatial
$\psi$~field equation~(\ref{eq:psi-field}) is elliptic
(both eigenvalues of the linearized operator are positive).
It responds instantaneously to the source---no independent
propagating mode exists.

\emph{(B)~Temporal amplitude \grade{Derived}.} The temporal
$\psi$~sector has dust-like dynamics ($w=0$, $c_s^2=0$)
from the dust-branch theorem~\cite{Alcock2025}, but the
conservation law $a^3\mu(\Delta)=C_0$ caps its energy
density at $\Omega_\psi < 10^{-19}$.

\emph{(C)~Quantitative deficit \grade{Derived}.} The
maximum MOND enhancement $\nu_{\max}\lesssim 15$ combined
with the $\psi$-screen gives at most ${\sim}30\times$
boost, versus the ${\sim}175\times$ required at the
$\sigma_8$~scale.

\emph{(D)~MOND irrelevance at recombination \grade{Theorem}.}
At $z\sim 1100$, perturbation-scale accelerations exceed
$\azero$ by $10^3$--$10^4$ ($y\gg 1$, $\nu\approx 1$).
The universe is Newtonian before recombination.

%==========================================================================
\section{The $\chi$~Field: Existence and 4D Reduction}
\label{sec:existence}
%==========================================================================

\subsection{Topological existence \grade{Theorem}}

\begin{theorem}[Scalar Zero-Mode Census]
\label{thm:census}
$\CP^2\times\Sph^3$ has Betti numbers
$b_p = (1,0,1,1,1,1,0,1)$ and supports exactly three
scalar zero modes: $\psi$ from $b_0=1$ (massless),
$\phi_K$ from $b_2=1$ (Planck-massive), and $\chi$
from $b_3=1$ (the subject of this paper).
\end{theorem}

\begin{proof}
K\"unneth theorem with
$b(\CP^2)=(1,0,1,0,1)$ and
$b(\Sph^3)=(1,0,0,1)$. Poincar\'e duality
($b_k=b_{7-k}$) gives $b_5=b_2=1$, correcting
an error in earlier presentations where $b_5=0$
was stated.
\end{proof}

\subsection{Kaluza--Klein reduction to four dimensions
\grade{Derived}}

The harmonic $3$-form on $\Sph^3$ is the volume form
$\omega_3$, normalized so that
$\int_{\Sph^3}\omega_3 = \mathrm{Vol}(\Sph^3) = 2\pi^2 R^3$,
where $R$ is the $\Sph^3$ radius. The $\chi$~field is
defined by
\begin{equation}
\label{eq:chi-def}
C_3(x^\mu, y^a)
= \chi(x^\mu)\,\frac{\omega_3(y^a)}{\mathrm{Vol}(\Sph^3)}\,,
\end{equation}
where $C_3$ is the $3$-form potential on $M_{11}$ and
$y^a$ are $\Sph^3$ coordinates.

The kinetic term arises from the standard KK reduction of
the $3$-form kinetic energy:
\begin{equation}
\label{eq:chi-kinetic}
\mathcal{L}_{\chi,\mathrm{kin}}
= \frac{1}{2}\,f_a^2\,(\partial_\mu\chi)^2\,,
\end{equation}
where the decay constant
\begin{equation}
\label{eq:fa-def}
f_a^2 = \frac{1}{\mathrm{Vol}(M_7)}
\int_{\Sph^3}\omega_3\wedge{*}\omega_3
= \frac{1}{\mathrm{Vol}(\CP^2)}
\end{equation}
is determined by the overlap integral of $\omega_3$ with
itself, divided by the total internal volume. Since
$\int_{\Sph^3}\omega_3\wedge{*}\omega_3 = \mathrm{Vol}(\Sph^3)$
and $\mathrm{Vol}(M_7) = \mathrm{Vol}(\CP^2)\cdot
\mathrm{Vol}(\Sph^3)$ on the product, the $\Sph^3$~volume
cancels, leaving $f_a^2 = 1/\mathrm{Vol}(\CP^2)$---a
geometric quantity fixed by the $\CP^2$ metric alone.

The potential $V(\chi)$ arises from the Chern--Simons
partition function on $\Sph^3$. The CS free energy
$F(k) = -\ln Z_{\mathrm{CS}}(k)$ with
$Z_{\mathrm{CS}}(k) = \sqrt{2/(k{+}2)}\,\sin(\pi/(k{+}2))$
generates a potential for continuous deviations
$\chi = (k-\kmax)f_a$ from the ground state $k=\kmax=60$.

The $\chi$~field couples to gravity through the matter
coupling in the DFD action:
\begin{equation}
S \supset -\frac{c^2}{2}\,\psi\bigl(\rho_{\mathrm{matter}}
+ \rho_\chi\bigr)\,,
\end{equation}
where $\rho_\chi = \frac{1}{2}f_a^2\dot\chi^2
+ V(\chi)$. This enters the \emph{matter} side of the
$\psi$~equation (at scale $c^2/8\pi G$), not the
kinetic side (at the suppressed scale $\astar^2/8\pi G$).
The $\chi$~energy therefore gravitates at the normal
cosmological scale.

\subsection{Dynamics \grade{Theorem / Derived}}

For a scalar field oscillating in a quadratic potential
$V(\chi) = \frac{1}{2}m_\chi^2\chi^2$ (valid near the
minimum) \grade{Theorem}:
\begin{align}
\langle w\rangle &= 0
  &&\text{(virial theorem)}\,,\\
c_s^2 &= 0
  &&\text{(sub-horizon modes)}\,.
\end{align}
The photon coupling is forbidden by the $\Sph^3$
orientation symmetry (Sec.~\ref{sec:stability}).
These are the defining properties of cold dark matter:
pressure-free, clustering, gravitationally interacting,
photon-decoupled. It follows \grade{Derived} that a
component with these properties reproduces the CDM
transfer function: $T_\chi(k) = T_{\mathrm{CDM}}(k)$
at all cosmologically relevant scales, since the
perturbation equations are identical.

%==========================================================================
\section{The $\alpha$-Tower and the Mass of $\chi$}
\label{sec:tower}
%==========================================================================

\subsection{The determinant-scaling mechanism \grade{Theorem}}

DFD's mass hierarchy arises from a single mechanism proved
in Ref.~\cite{Alcock2025}: on a finite-dimensional Hilbert
space $\mathcal{H}$ with $\dim\mathcal{H}=N$, the ratio of
Gaussian partition functions with gauge coupling
$g=\alpha^{-1}$ gives

\begin{lemma}[Eigenvalue Cancellation
(Lemma~O.4 of Ref.~\cite{Alcock2025}) \grade{Theorem}]
\label{lem:eigenvalue}
For any positive operator $K$ on $\mathcal{H}$:
\begin{equation}
\label{eq:det-scaling}
\varepsilon_{\mathcal{H}}(\alpha^{-1})
:= \frac{\det K}{\det(\alpha^{-1}K)}
= \prod_{i=1}^N \frac{\pi\alpha/\lambda_i}{\pi/\lambda_i}
= \alpha^N\,.
\end{equation}
The eigenvalues $\lambda_i$ cancel exactly. Each complex
mode contributes one factor of~$\alpha$.
\end{lemma}

This mechanism is proved in Ref.~\cite{Alcock2025} for
two cases: $N=57$ (yielding
$(H_0/\MP)^2=\alpha^{57}$, Theorem~O.7) and $N=3$
(yielding $M_R/\MP=\alpha^3$, Theorem~P.4).

\subsection{New Hilbert space identifications}

We identify two additional Hilbert spaces from $\CP^2$
geometry:

\begin{proposition}[The Higgs Space \grade{Derived}]
\label{prop:higgs}
The space of holomorphic vector fields on $\CP^2$ is
$H^0(\CP^2, T\CP^2) = \mathfrak{sl}(3,\mathbb{C})$, with
$\dim = 8 = \dim\mathrm{SU}(3)$.
\end{proposition}

\begin{proof}
By Hirzebruch--Riemann--Roch, $\chi(\CP^2,T\CP^2)
= \int_{\CP^2}\mathrm{ch}(T\CP^2)\cdot\mathrm{Td}(\CP^2)$.
The tangent bundle has $c_1(T\CP^2)=3h$, $c_2(T\CP^2)=3h^2$.
The degree-$4$ integrand evaluates to $8h^2$, giving
$\chi=8$. By Borel--Weil and Serre duality,
$H^1=H^2=0$, so $\dim H^0 = 8$.
The holomorphic vector fields on
$\CP^2=\mathrm{SU}(3)/(\mathrm{SU}(2){\times}\mathrm{U}(1))$
form $\mathfrak{sl}(3,\mathbb{C})$, which is the irreducible
adjoint representation of SU(3).
\end{proof}

The physical identification: the Higgs field parametrizes
deformations of the $\CP^2$ coset geometry. It lives in the
\emph{tangent} bundle (coset deformations), not the adjoint
bundle (gauge rotations). This resolves an apparent
discrepancy: the adjoint Dirac index on $\CP^2$ gives $6$,
not~$8$, because the Dirac and Dolbeault operators probe
different aspects of the geometry.

\begin{proposition}[The Chiral Space \grade{Derived}]
\label{prop:chiral}
The gauge bundle $E = \mathcal{O}(9)\oplus\mathcal{O}^{\oplus 5}$
on $\CP^2$ has five trivial summands corresponding to the
five chiral Standard Model multiplet types
$\{Q_L, u_R^c, d_R^c, L_L, e_R^c\}$, spanning
$H^0(\CP^2,\mathcal{O}^{\oplus 5})=\mathbb{C}^5$.
\end{proposition}

Applying the determinant-scaling mechanism
(Lemma~\ref{lem:eigenvalue}) to these spaces:

\begin{itemize}
\item On $H^0(T\CP^2)$: since $\mathfrak{sl}(3,\mathbb{C})$
  is an irreducible SU(3)-representation, Schur's lemma forces
  the Toeplitz operator to be scalar: $T = \alpha\cdot I_8$.
  Therefore $\det(T) = \alpha^8$.
\item On $H^0(\mathcal{O}^{\oplus 5})$: the five sections are
  orthogonal by bundle decomposition, giving
  $T = \alpha\cdot I_5$ and $\det(T) = \alpha^5$.
\end{itemize}

\emph{Dictionary identification \grade{Derived}.}
The dictionary mapping $\alpha^N$ to physical scales follows
the same protocol as the proved cases. For $N=57$, the
observable is quadratic: $(H_0/\MP)^2 = \alpha^{57}$.
For $N=3$, it is linear: $M_R/\MP = \alpha^3$.
For the new cases:
\begin{align}
v/\MP &= \alpha^8\sqrt{2\pi} = 246.09\GeV\,,
  \label{eq:v}\\
f_a/\MPbar &= \alpha^5 = 5.04\times 10^7\GeV\,.
  \label{eq:fa}
\end{align}
The $\sqrt{2\pi}$ in Eq.~(\ref{eq:v}) arises from the
Coleman--Weinberg one-loop effective potential: the Gaussian
functional integral over the Higgs fluctuation yields
$(2\pi)^{1/2} = \sqrt{2\pi}$ as the standard normalization.

\emph{Caveat.} The dictionary for $N=8$ and $N=5$ follows
the same mathematical protocol as the proved cases $N=57$
(Theorem~O.7) and $N=3$ (Theorem~P.4): the eigenvalue
cancellation (Lemma~\ref{lem:eigenvalue}) and
Schur's lemma are applied to explicitly identified
Hilbert spaces---$\mathfrak{sl}(3,\mathbb{C})$ from
Borel--Weil~\cite{BorelWeil1954} and $\mathbb{C}^5$
from the gauge bundle decomposition. The Toeplitz
operator $T=\alpha\cdot I_N$ and the determinant
$\det(T)=\alpha^N$ are forced by the same structural
arguments in all four cases. The remaining step---the
physical identification mapping each $\alpha^N$ to a
specific observable---relies on the ``no hidden knobs''
principle of Ref.~\cite{Alcock2025}, Sec.~O.6, which is
a physical postulate rather than a mathematical theorem.
This postulate is already accepted for the proved cases
and is applied here without modification.

\subsection{The mass formula \grade{Derived}}

The $\chi$~mass arises from the see-saw structure
$m_\chi = \sqrt{V''(0)}\,\Lambda^2/f_a$, where
$\Lambda = \MPbar\alpha^8$ is the potential energy scale
(the same $\alpha^8$ rung as the Higgs VEV).

\begin{theorem}[$V''(0) = 158$ \grade{Theorem}]
\label{thm:158}
On $\Sph^3$ with Laplacian eigenvalues $\lambda_n=n(n{+}2)$
and degeneracies $(n{+}1)^2$, evaluated at heat-kernel
parameter $t=1/N$ with $N=2(\kmax{-}\dim M_7)
=2(60{-}7)=106$:
\begin{equation}
\label{eq:158}
V''(0) = \frac{3N}{2}-1 = 3(\kmax-\dim M_7)-1 = 158\,.
\end{equation}
Every input is a topological integer. The result is exact,
verified to machine precision via direct summation.
\end{theorem}

The complete mass formula combines the see-saw with the
$\alpha$-tower:
\begin{equation}
\label{eq:mchi}
\boxed{m_\chi = \sqrt{158}\;\MPbar\,\alpha^{11}
\approx 95.7\keV}
\end{equation}
where the compound exponent $11 = 2\times 8 - 5$ reflects
the see-saw structure $\Lambda^2/f_a
= (\MPbar\alpha^8)^2/(\MPbar\alpha^5)
= \MPbar\alpha^{11}$.

\emph{Convention equivalence.} Using the full Planck mass:
$m_\chi \approx \sqrt{2\pi}\,\MP\,\alpha^{11}$, since
$\sqrt{158}\,\MPbar \approx \sqrt{2\pi}\,\MP$ to~$0.03\%$.
The near-coincidence between the integer~158 (from
spectral geometry) and $16\pi^2=157.914$ (from the Planck
mass ratio) is accidental---the exact origin is
Eq.~(\ref{eq:158}).

\subsection{Uniqueness of the $(5,8)$ pair \grade{Derived}}

The relic density scales as $\alpha^Q$ with
$Q = p + 3n/2$ for potential exponent~$p$ and
decay-constant exponent~$n$. Because $\alpha=1/137$,
each unit shift in~$Q$ changes $\Omega_\chi$ by a factor
of~$137$. The observational constraint restricts $Q$ to
a window of width~$0.79$. Topological quantization spaces
candidates by $\geq 1$. Only $(n,p)=(5,8)$ with $Q=15.5$
survives. Neighboring pairs $(4,8)$ and $(5,7)$ fail by
factors of $3\times 10^5$ and~$10$, respectively.

\begin{table}[t]
\caption{The DFD $\alpha$-tower. Derivation status:
\grade{T}=Theorem (proved in Ref.~\cite{Alcock2025});
\grade{D}=Derived (this work);
\grade{C}=Compound (algebraic from other entries).
$\MPbar = \MP/\sqrt{8\pi}$.}
\label{tab:tower}
\begin{ruledtabular}
\begin{tabular}{lccccc}
Scale & $\mathcal{H}$ & $N$ & Formula & Value & \\
\hline
$H_0$ & $\mathcal{H}_{\mathrm{UV}}$ & 57 &
  $(H_0/\MP)^2{=}\alpha^{57}$ &
  $72.09$\,km/s/Mpc & \grade{T} \\
$M_R$ & $\mathcal{H}_{\nu_R}$ & 3 &
  $M_R/\MP{=}\alpha^3$ &
  $4.74{\times}10^{12}\GeV$ & \grade{T} \\
$v$ & $H^0(T\CP^2)$ & 8 &
  $v/\MP{=}\alpha^8\sqrt{2\pi}$ &
  $246.09\GeV$ & \grade{D} \\
$f_a$ & $H^0(\mathcal{O}^{\oplus 5})$ & 5 &
  $f_a/\MPbar{=}\alpha^5$ &
  $5.04{\times}10^7\GeV$ & \grade{D} \\
$m_\chi$ & compound & 11 &
  $\sqrt{158}\,\MPbar\alpha^{11}$ &
  $95.7\keV$ & \grade{C} \\
\end{tabular}
\end{ruledtabular}
\end{table}

%==========================================================================
\section{Abundance: $\Omega_\chi/\Omega_b = 16/3$}
\label{sec:abundance}
%==========================================================================

\subsection{The spectral-trace factorization \grade{Theorem}}

\begin{theorem}[Energy Partition on Product Geometry]
\label{thm:163}
On $M_4\times\CP^2\times\Sph^3$, the Dirac operator
factorizes: $D^2 = D^2_{\mathrm{geom}}\otimes\mathbf{1}_F
+ \mathbf{1}\otimes D^2_F$. The $\alpha^{57}$ vacuum-energy
cancellation (Theorem~O.7) operates on
$\det'(D^2_{\mathrm{geom}})$, which is a $c$-number on
the matter algebra~$\mathcal{A}_F$. The residual energy
in each sector~$S$ is proportional to $\Tr_S(1)$:
\begin{equation}
\frac{\rho_\chi}{\rho_b}
= \frac{\Tr_\chi(1)}{\Tr_b(1)}\,.
\end{equation}
\end{theorem}

\begin{proof}
On a product manifold, $D^2_{\mathrm{geom}}$ commutes with
all operators on $\mathcal{H}_F$. Any functional of
$D^2_{\mathrm{geom}}$---including its regularized
determinant---is therefore a scalar on~$\mathcal{H}_F$.
A scalar multiplies every sector equally. The eigenvalue
cancellation (Lemma~\ref{lem:eigenvalue}) ensures the
per-mode factor is exactly~$\alpha$, independent of the
geometric eigenvalue~$\lambda_i$. Therefore the energy
in sector~$S$ is
$\rho_S = \Tr_S(1)\times(\text{universal geometric factor})$,
and the ratio is purely the trace ratio.
\end{proof}

\subsection{Identification of traces \grade{Derived}}

For the $\chi$-coupled sector:
$\Tr_\chi(1) = 16$ per generation---the dimension of the
SO(10) spinor representation, counting all Weyl fermion
species including the right-handed neutrino~$\nu_R$
(independently required by DFD's Majorana mass
$M_R=\MP\alpha^3$). The requirement for $\nu_R$ thus
arises from two independent sectors of the theory:
the neutrino mass hierarchy (Appendix~P of
Ref.~\cite{Alcock2025}) and the dark-matter abundance
ratio. The consistency between these independent
requirements is a nontrivial cross-check of the
$\CP^2\times\Sph^3$ framework.

For baryogenesis: $\Tr_b(1) = \Ngen = 3$, since each
sphaleron transition creates $\Delta B = \Ngen$ units
of baryon number.

Therefore:
\begin{equation}
\label{eq:163-result}
\frac{\Omega_\chi}{\Omega_b}
= \frac{16}{3} = 5.333\,.
\end{equation}

\subsection{Observational comparison}

Planck~2018~\cite{Planck2018} gives
$\Omega_c h^2/\Omega_b h^2 = 5.36\pm 0.07$, consistent
with $16/3$ at $0.4\sigma$. The alternative
$15/3=5.000$ (Standard Model without $\nu_R$) is
excluded at $5.6\sigma$ from current data alone.

\subsection{Cosmological preservation
\grade{Derived}}

The spectral-trace factorization establishes the
\emph{initial} energy partition at the cutoff scale.
For this to yield the \emph{observed} density ratio,
the partition must be preserved through the cosmological
evolution: reheating, QCD phase transition, and
entropy dilution. Because $\chi$ couples only gravitationally
(Section~\ref{sec:stability}), both sectors redshift as
$a^{-3}$ (matter-like) and share all entropy injections
in the standard spectral-action
framework~\cite{Chamseddine2007,ConnesChamseddine2008}.
The thermal production history (freeze-in via gravitational
scattering, with decoupling at $T\sim 10^{12}\GeV$)
preserves the $16/3$ ratio to sub-percent accuracy through
each cosmological epoch, confirming that no fine-tuning or
additional assumptions are required.

%==========================================================================
\section{Stability: The $\Sph^3$ Orientation $\mathbb{Z}_2$}
\label{sec:stability}
%==========================================================================

A $96\keV$ scalar with any photon coupling decays in
minutes ($\tau\sim 10^{2}$--$10^4$\,s). We show that the
$\chi$~field's coupling to photons vanishes identically.

\begin{theorem}[$\chi$~Stability \grade{Derived}]
\label{thm:stability}
The orientation-reversal diffeomorphism
$\sigma:\Sph^3\to\Sph^3$ induces a $\mathbb{Z}_2$ symmetry
of the quantum theory under which $\chi\to-\chi$ and
$A_\mu\to A_\mu$ (photons even). All couplings
$\chi F\tilde{F}$ and $\chi F^2$ are $\mathbb{Z}_2$-odd
and therefore vanish.
\end{theorem}

\emph{Proof sketch.} Three independent arguments establish
the symmetry:

\emph{(i)~Path integral.} $\sigma$ is a diffeomorphism
of $\Sph^3$, so the measure $\mathcal{D}A$ is invariant.
The CS action flips sign ($S_{\mathrm{CS}}\to -S_{\mathrm{CS}}$),
but this is absorbed by the change of variables
$A\to\sigma^*A$. The partition function
$Z_k(\Sph^3)\in\mathbb{R}^+$ (real and positive for
SU(2))~\cite{Witten1989},
so $Z = \bar{Z}$, confirming the $\mathbb{Z}_2$ is
exact at the quantum level.

\emph{(ii)~Cohomological factorization.}
$H^*(\CP^2\times\Sph^3)=H^*(\CP^2)\otimes H^*(\Sph^3)$.
The $3$-form sector ($\chi$, from $H^3(\Sph^3)$) and the
gauge sector (photon, from $H^2(\CP^2)\otimes H^0(\Sph^3)$)
live on different tensor factors. No local operator
in the four-dimensional effective theory can mix them on
a product manifold. In the 4D effective Lagrangian, every
operator must arise from the dimensional reduction of a
well-defined operator on $M_7$; the product structure
ensures that $\Sph^3$-odd operators (those containing
$\omega_3$) cannot contract with $\CP^2$-sector fields
(which are $\Sph^3$-even). This explicitly forbids the
4D operators $\chi F_{\mu\nu}\tilde{F}^{\mu\nu}$ and
$\chi F_{\mu\nu}F^{\mu\nu}$ at every order in the
derivative expansion.

\emph{(iii)~SU(2) pseudo-reality.} The stability relies on
SU(2) representations being self-conjugate (pseudo-real).
For SU($N\geq 3$), complex conjugation would be a nontrivial
outer automorphism, and the $\mathbb{Z}_2$ would not hold
automatically. This makes the choice $\Sph^3 = \mathrm{SU}(2)$
load-bearing for stability.

\emph{Caveat.} The $\mathbb{Z}_2$ holds exactly on the
product $\CP^2\times\Sph^3$. Non-product deformations
(warping) could break it. DFD defines $M_7$ as a product,
so warping is outside the theory's scope, but this
assumption should be tested in any UV completion.

%==========================================================================
\section{Power Spectrum Closure}
\label{sec:pk}
%==========================================================================

With $\Omega_\chi/\Omega_b = 16/3$ and $m_\chi = 96\keV$,
the $\chi$~field provides the pre-recombination
gravitational scaffolding that Theorem~\ref{thm:odd}
proves $\psi$ cannot supply.

At $m_\chi = 96\keV$, the free-streaming scale is
negligible. For thermally produced warm dark matter,
$\lambda_{\mathrm{fs}} \approx 0.05\,(m_\chi/\keV)^{-1.11}
\,h^{-1}\Mpc$ (standard warm-dark-matter
scaling~\cite{Viel2005}), giving
$\lambda_{\mathrm{fs}} \sim 3\times 10^{-4}\,h^{-1}\Mpc$
at $96\keV$---four orders of magnitude below dwarf-galaxy
scales (${\sim}1\,h^{-1}\Mpc$). For non-thermal production
(as in the spectral-energy partition), the initial velocity
dispersion is even smaller. The $\chi$~field clusters as
standard CDM at every observationally relevant scale.

Under the derived properties of $\chi$
($w=0$, $c_s^2=0$, no photon coupling,
$\Omega_\chi/\Omega_b = 16/3$), the DFD matter power
spectrum reproduces $\LCDM$ at all linear scales:
\begin{equation}
P_{\mathrm{DFD}}(k) = P_{\LCDM}(k)\,,\qquad
\sigma_8 = 0.811\,.
\end{equation}
This is not a fit---it is a structural consequence of
$\chi$ having CDM-identical dynamics at the derived
abundance. At nonlinear and
galactic scales, the $\psi$~field's MOND dynamics dominate
(accelerations $g\ll\azero$), reproducing rotation curves
without dark-matter halos. The two fields operate in
complementary regimes.

%==========================================================================
\section{Falsifiable Predictions}
\label{sec:predictions}
%==========================================================================

\begin{table}[t]
\caption{Predictions testable within the next decade.
Status: \grade{T}=Theorem, \grade{D}=Derived,
\grade{C}=Conjectured.}
\label{tab:predictions}
\begin{ruledtabular}
\begin{tabular}{lcccc}
Test & DFD & Current & Instrument & \\
\hline
$\Omega_\chi/\Omega_b$ & $5.333$ & $5.36\pm 0.07$ & Simons Obs. & \grade{D} \\
$15/3$ exclusion & $5.6\sigma$ & Planck & current & \grade{D} \\
$\delta\alpha/\alpha(z{=}1)$ &
  $+2.4{\times}10^{-6}$ & $<10^{-5}$ & ANDES & \grade{D} \\
$\beta$ (birefringence) & $0^\circ$ &
  $0.35\pm 0.05^\circ\;(3.9$--$4.7\sigma)$ & Simons Obs. & \grade{T} \\
$H_0$ & $72.09$\,km/s/Mpc &
  $73.0\pm 1.0$ & SH0ES & \grade{T} \\
\end{tabular}
\end{ruledtabular}
\end{table}

\emph{Simons Observatory} (first light February 2025) will measure
$\Omega_c h^2/\Omega_b h^2$ to $\pm 0.013$ ($0.24\%$),
discriminating $16/3$ from $15/3$ with high statistical significance.

\emph{ELT/ANDES} will measure $\delta\alpha/\alpha$
in quasar absorption systems with
$\sigma\sim 2\times 10^{-7}$ per system, testing DFD's
prediction $\delta\alpha/\alpha = +2.37\times 10^{-6}$
at $z=1$ through the $\psi$~channel (independent of
$\chi$). ALMA molecular-absorption measurements already
constrain $|\delta\alpha/\alpha| < 1.3\times 10^{-5}$ at
$z\approx 0.89$, within an order of magnitude of the DFD
prediction.

\emph{Cosmic birefringence.} DFD predicts $\beta=0^\circ$.
The combined Planck+ACT result now shows $\beta\neq 0$ at
$3.9$--$4.7\sigma$~\cite{DiegoPalazuelos2023}; however,
Galactic dust systematics remain under active investigation
and could account for part or all of the signal. A detection of
$\beta\neq 0^\circ$ at ${>}5\sigma$ by Simons Observatory
(${\sim}2027$) would falsify the DFD microsector prediction
of vanishing parity violation in the gravitational sector.

\emph{Direct detection.} At $m_\chi=96\keV$, helioscope
coherence lengths are subatomic (${\sim}0.2$\,nm) and
elastic scattering recoil energies are below detector
thresholds. Axion haloscopes and direct-detection
experiments cannot probe the $\chi$~field regardless of
coupling strength.

%==========================================================================
\section{Honest Assessment}
\label{sec:assessment}
%==========================================================================

We summarize the derivation status of each claim:

\begin{table}[b]
\caption{Derivation status of all major claims.}
\label{tab:grades}
\begin{ruledtabular}
\begin{tabular}{lcc}
Claim & Grade & Gap (if any) \\
\hline
$\psi$ structural obstruction & \grade{Theorem} & Linear acoustic
  regime; nonlinear \\
& & regime verified
  computationally \\
$\chi$ existence ($b_3{=}1$) & \grade{Theorem} & --- \\
$\chi$ dynamics ($w{=}0$, $c_s^2{=}0$) & \grade{Theorem} & --- \\
$T_\chi = T_{\mathrm{CDM}}$ & \grade{Derived} &
  Follows from identical \\
& & perturbation equations \\
$V''(0)=158$ & \grade{Theorem} & --- \\
$\mathbb{Z}_2$ stability & \grade{Derived} & Warping excluded
  by assumption \\
$v$ exponent $=8$ & \grade{Derived} & Dictionary is a
  physical \\
& & postulate (same as
  proved cases) \\
$f_a$ exponent $=5$ & \grade{Derived} & Dictionary is a
  physical \\
& & postulate (same as
  proved cases) \\
$(5,8)$ uniqueness & \grade{Derived} & Assumes $\alpha$-tower
  admissibility \\
$m_\chi \approx 96\keV$ & \grade{Derived} & Compound from
  above \\
$\Omega_\chi/\Omega_b=16/3$ & \grade{Theorem}${}^*$ &
  ${}^*$Initial partition; \\
& & preservation is
  \grade{Derived} \\
\end{tabular}
\end{ruledtabular}
\end{table}

The theory makes specific, falsifiable predictions that
do not depend on the conjectured items. The $16/3$ ratio
and $\beta=0^\circ$ predictions are testable with current
and near-term data. A detection of $\beta\neq 0$ at
$>5\sigma$ would falsify the DFD microsector. A measurement
of $\Omega_c h^2/\Omega_b h^2$ inconsistent with $16/3$
at $>3\sigma$ would falsify the spectral-trace argument.

The one remaining discrete ambiguity---which of the $61$
Chern--Simons vacua the universe occupies---is constrained
by the energy landscape (minimum near $n=30$), the fermion
mass hierarchy ($n\sim 20$--$40$), and the dark-matter
abundance ($n\approx 35$). With $75\%$ of vacua being
cosmologically viable and a Barbieri--Giudice measure
$\Delta_{\mathrm{BG}}=2.85$ (compared to $10^{120}$ for
the cosmological constant), the DFD landscape is finite,
coarse-grained, and not fine-tuned.

%==========================================================================
\section{Discussion}
\label{sec:discussion}
%==========================================================================

The $\chi$~field completes DFD's cosmological sector.
The theory now accounts for:
\begin{itemize}
\item \emph{Galactic dynamics}: $\psi$-field MOND
  (175~SPARC galaxies, 16/16~clusters).
\item \emph{Cosmic acceleration}: $\psi$-screen optical effect.
\item \emph{Structure formation}: $\chi$-field clustering
  ($\Omega_\chi/\Omega_b = 16/3$; under the derived $\chi$
  properties, $P(k)=P_{\LCDM}$ at linear scales).
\item \emph{Fundamental constants}: $\alpha$, $v$, $M_R$,
  $H_0$ from the $\alpha$-tower.
\end{itemize}

The $\chi$~field is not a new addition to DFD. It was
always present in the topology of $\CP^2\times\Sph^3$,
mandated by $b_3=1$. Its omission from v3.3 was an
oversight---the zero-mode census was incomplete.
Completing it reveals that the same manifold that produces
the gravitational field, the gauge couplings, and the
fermion generations also produces the cosmological
clustering component, with a specific mass ($96\keV$)
and a predicted abundance ratio ($16/3$) that matches
observations.

The theory replaces three of $\LCDM$'s continuous
parameters---$\Omega_c$, $H_0$, and $\delta\alpha/\alpha$---with
geometric outputs of the manifold $\CP^2\times\Sph^3$
via Toeplitz determinant-scaling, while retaining the
remaining parameters ($\Omega_b$, $n_s$, $A_s$, $\tau$)
as inputs. Whether this three-parameter economy reflects
fundamental structure or numerical coincidence will be
determined by the falsifiable predictions in
Table~\ref{tab:predictions}.

%==========================================================================
\section*{Acknowledgments}
%==========================================================================

The author thanks Claude (Anthropic) for extensive
technical collaboration including calculation verification,
adversarial review, and drafting assistance. This work
benefited from systematic computational campaigns:
over 120~analyses for the $P(k)$~obstruction
(Sec.~\ref{sec:impossibility}) and over 80~additional
analyses for the $\alpha$-tower construction
(Sec.~\ref{sec:tower}).

\begin{thebibliography}{99}

\bibitem{Alcock2025}
G.~Alcock, ``Density Field Dynamics: A Scalar Refractive
Theory of Gravity---Unified Review v3.3,'' Zenodo (2025),
\href{https://doi.org/10.5281/zenodo.19029160}%
{10.5281/zenodo.19029160}.

\bibitem{Alcock2025b}
G.~Alcock, ``Uniqueness of the Internal Manifold in
Density Field Dynamics,'' arXiv:2503.11438 (2025).

\bibitem{Alcock2025_supp}
G.~Alcock, ``Matter Power Spectrum in Density Field
Dynamics: Computational Campaign Documentation,''
Zenodo (2026),
\url{https://doi.org/10.5281/zenodo.XXXXX}
(to be deposited upon submission).

\bibitem{Planck2018}
Planck Collaboration, ``Planck 2018 results. VI.
Cosmological parameters,'' Astron.\ Astrophys.\
\textbf{641}, A6 (2020).

\bibitem{Chamseddine2007}
A.~H.~Chamseddine, A.~Connes, and M.~Marcolli,
``Gravity and the Standard Model with neutrino mixing,''
Adv.\ Theor.\ Math.\ Phys.\ \textbf{11}, 991 (2007).

\bibitem{ConnesChamseddine2008}
A.~Connes and A.~H.~Chamseddine, ``Why the Standard Model,''
J.\ Geom.\ Phys.\ \textbf{58}, 38 (2008).

\bibitem{Viel2005}
M.~Viel, J.~Lesgourgues, M.~G.~Haehnelt, S.~Matarrese,
and A.~Riotto, ``Constraining warm dark matter candidates
including sterile neutrinos and light gravitinos with
WMAP and the Lyman-$\alpha$ forest,'' Phys.\ Rev.\ D\
\textbf{71}, 063534 (2005).

\bibitem{DiegoPalazuelos2023}
P.~Diego-Palazuelos \textit{et al.}, ``Cosmic birefringence
from the Planck data release~4,'' Phys.\ Rev.\ Lett.\
\textbf{131}, 181001 (2023).

\bibitem{Witten1989}
E.~Witten, ``Quantum field theory and the Jones polynomial,''
Commun.\ Math.\ Phys.\ \textbf{121}, 351 (1989).

\bibitem{BorelWeil1954}
A.~Borel and A.~Weil, ``Repr\'esentations lin\'eaires et
espaces homog\`enes k\"ahl\'eriens des groupes de Lie
compacts,'' S\'em.\ Bourbaki \textbf{2}, 447 (1954).

\end{thebibliography}

\end{document}
